\begin{titlepage}
    \pagenumbering{gobble}
    \centering

    % Logo institucional UNJBG
    \begin{tikzpicture}
        % Escudo circular estilizado UNJBG
        \draw[line width=3pt, color=mdpiBlue] (0,0) circle (1.5cm);
        \draw[line width=1.5pt, color=mdpiBlue] (0,0) circle (1.1cm);
        \node[align=center, text=mdpiBlue, font=\bfseries\scriptsize] at (0, 0.55) {U N J B G};
        \node[align=center, text=mdpiBlue, font=\tiny] at (0, 0.1) {TACNA -- PERÚ};
        \draw[color=mdpiBlue, line width=1pt] (-0.7,-0.25) -- (0.7,-0.25);
        \node[align=center, text=mdpiBlue, font=\bfseries\large] at (0, -0.7) {\faUniversity};
        \draw[color=mdpiBlue!60, line width=0.5pt] (0,0) circle (1.3cm);
    \end{tikzpicture}
    \par\vspace{0.2cm}
    {\fontsize{13}{15}\bfseries\color{mdpiBlue} Universidad Nacional Jorge Basadre Grohmann\par}
    {\fontsize{12}{14}\selectfont Escuela Profesional de Ingeniería en Informática y Sistemas\par}

    \vspace{0.6cm}
    {\fontsize{14}{16}\bfseries REDES I\par}
    {\fontsize{12}{14}\bfseries MEMORIA TÉCNICA\par}
    \vspace{0.2cm}
    {\fontsize{12}{14}\bfseries Diseño de la Infraestructura Física y Cableado Estructurado\\
    de una Red LAN — Policlínico de Especialidades Médicas\\
    ``Salud \& Vida'' S.A.C.\par}

    \vspace{0.5cm}

    \begin{center}
        {\bfseries Docente:}\\[0.2cm]
        \normalsize Mg. REYNALDO ALFONTE ZAPANA
    \end{center}

    \vspace{0.4cm}

    \begin{flushleft}
        {\bfseries Autores:}\\[0.4cm]
        \small
        \autorconorcid{Cachicatari Quito Aron Rene}{2024-119011}{0009-0001-0343-6817}\\[0.3cm]
        \autorsinorcid{Banchello Pino Rojas Raul Farid}{2024-119009}\\[0.3cm]
        \autorconorcid{Nina Estuco William Andre}{2024-119006}{0009-0002-6296-4473}\\[0.3cm]
        \autorsinorcid{Mamani Lima Josue Fernando}{2025-19055}\\[0.3cm]
        \autorconorcid{Briceño Pihuaycho Fernando Jesus}{2024-119029}{0009-0004-2188-6274}\\[0.3cm]
    \end{flushleft}

    \vspace{0.3cm}

    \begin{center}
        \begin{minipage}{0.96\textwidth}
            {\small \textbf{RESUMEN:} La presente memoria técnica documenta el diseño completo de la infraestructura de red LAN del Policlínico de Especialidades Médicas ``Salud \& Vida'' S.A.C., una edificación de 450~m$^2$ de planta única. El estudio cubre la definición de requerimientos (48~tomas activas + 8~de reserva = 56~puertos totales), la selección de topología física en estrella jerárquica con medios UTP~Cat~6 LSZH y F/UTP~Cat~6A blindado, la ingeniería de cableado estructurado bajo normas ANSI/TIA-568, 569, 606 y 607, el diseño físico del rack 42U con el cuadro de materiales valorizado en USD~\$9,632 y el protocolo de construcción y certificación del cableado bajo el estándar T568B, con distancia máxima de enlace permanente de 34.8~m, en un amplio margen sobre el límite normativo de 90~m.\par}

            \vspace{0.2cm}
            {\small \textbf{Palabras clave:} Cableado estructurado, ANSI/TIA-568, Red LAN hospitalaria, Cat~6 LSZH, F/UTP Cat~6A, Rack 42U, UPS, Wi-Fi~6, Certificación de enlace.\par}

            \vspace{0.3cm}
            {\small \textbf{ABSTRACT:} This technical report documents the complete LAN network infrastructure design for the ``Salud \& Vida'' Medical Specialty Polyclinic, a single-floor 450~m$^2$ building. The study covers requirements definition (48 active outlets + 8 spare = 56 total ports), hierarchical star topology selection with UTP~Cat~6 LSZH and shielded F/UTP~Cat~6A media, structured cabling engineering under ANSI/TIA-568, 569, 606 and 607 standards, a 42U rack physical design with a material list valued at USD~\$9,632, and the cable construction and certification protocol under the T568B standard, with a maximum permanent link distance of 34.8~m — well within the 90~m normative limit.\par}

            \vspace{0.2cm}
            {\small \textbf{Keywords:} Structured cabling, ANSI/TIA-568, Hospital LAN, Cat~6 LSZH, F/UTP~Cat~6A, 42U rack, UPS, Wi-Fi~6, Link certification.\par}
        \end{minipage}
    \end{center}

    \vfill
    {\normalsize Tacna, Octubre 2026\par}
\end{titlepage}
